The study of the Bible has generated one of the largest and most enduring bodies of scholarly commentary in human history. Prominent theologians, from John Calvin in the sixteenth century to Archibald T. Robertson in the twentieth, have produced detailed analyses that illuminate the linguistic, theological, and spiritual dimensions of the biblical text. For many Bible studiers, these commentaries are indispensable. Yet the sheer volume of material presents a practical challenge: consulting even a handful of commentators for a single verse requires navigating multiple lengthy texts, synthesizing sometimes divergent perspectives, and extracting the most relevant insights.

BibleGoAI grew out of a simple goal: \textit{help people study the Bible using time-tested commentaries presented in a simple and accessible way}, not only for theological understanding but also to help believers in their spiritual growth.

Modern advances in natural language processing (NLP) and large language models (LLMs) offer new possibilities for addressing this challenge. Systems based on retrieval-augmented generation (RAG) have demonstrated the ability to ground LLM outputs in authoritative source documents \citep{lewis2020rag}, while sentence embeddings \citep{cer2018use, reimers2019sentencebert} enable fine-grained semantic comparison between generated and source texts. These capabilities make it feasible to build systems that not only summarize but also \textit{validate} the fidelity of their outputs against the original material.

In this paper, we present \textbf{BibleGoAI}, a multi-stage pipeline that transforms curated biblical commentary into concise, source-adherent summary paragraphs called \textit{Brief Helpful Texts} (BHTs). The name reflects the project's core value: bringing the richness of difficult-to-read commentary to the fingertips of Bible readers in an accessible way. The system is designed around three principles:

\begin{enumerate}[nosep,leftmargin=*]
    \item \textbf{Fidelity}: BHTs must be grounded in the actual words of trusted commentators, not in the LLM's own knowledge. The pipeline enforces this through hard word-proportion gates and prompt-level targets; see Section~\ref{sec:evaluation} for realized statistics.
    \item \textbf{Conciseness}: Each BHT is a short paragraph of roughly 80 words and a few sentences (typically 2--4; see Section~\ref{sec:evaluation} for per-pipeline distributions), making it practical for quick reference during Bible study.
    \item \textbf{Traceability}: Every claim in a BHT is linked back to its source commentary via footnotes with similarity scores, enabling readers to verify and explore further.
\end{enumerate}

The pipeline was run once over the entire Bible: the full 27-book New Testament (NT; 7{,}957 verses) and the full 39-book Old Testament (OT; 23{,}145 verses). The OT was served through two complementary pipeline variants: a single-pass regular pipeline and a two-stage split pipeline (see Section~\ref{sec:evaluation} for details). It draws from a curated corpus of 17 unique commentators (nine NT, twelve OT, four shared) spanning four centuries and organized into three tiers. The resulting BHTs are served to end users through BibleGo.org, the project's public web application. It is used by a small but active Bible-study community of roughly 100--300 unique users per month, with monthly page-view volumes ranging from ${\sim}200$ to ${\sim}2{,}500$ depending on seasonal and community activity.

Our contributions are as follows:
\begin{itemize}[nosep,leftmargin=*]
    \item A multi-stage pipeline for commentary synthesis combining LLM-based quote extraction, tiered prompt engineering, and constraint-enforced generation.
    \item A composite quality scoring framework that integrates semantic similarity and source fidelity checks.
    \item An OT-specific design that uses retrieval-augmented generation to narrow multi-verse commentary, served through two pipeline variants (single-pass regular and two-stage split) covering all 39 OT books, with per-pipeline metrics reported.
    \item A user-facing application (BibleGo.org) serving the generated BHTs to a live Bible-study community, demonstrating the practical impact of the system.
\end{itemize}

\paragraph{Project context.} BibleGoAI is a volunteer-driven collaboration among Christians, mostly based in the Cleveland, Ohio area, with backgrounds in software engineering and church ministry. The project proceeded through weekly round-table meetings in 2023, a volunteer commentary digitization effort (Henry Alford's 2{,}544-page \textit{New Testament for English Readers} commentary, digitized by 28 volunteers in 2023, and Robert Govett's 627-page \textit{Commentary on Revelation}, digitized through a similar process and completed in 2025), and an intensive April--May 2024 development sprint. The NT was publicly launched at a New Year's church conference in December 2023 and the OT at a Memorial Day church conference in 2024. Several technical decisions were shaped by this community context: the commentator corpus was curated by ministry leaders for general helpfulness in Bible study; usability testing was conducted with members of the project's Bible study community; the NASB~1995 translation was adopted after a request to use the ESV was declined by its publisher's licensing terms; NASB also provided the Strong's concordance tagging that backs the application's original-language word study; and a sustained cadence of internal demos served as the primary feedback mechanism.
